% =====================================================================
% Section 2 — Background
% =====================================================================

\section{Background}\label{sec:background}

\subsection{Valve-train failure modes at high rotational speed}
\label{sec:bg-failure-modes}

At elevated rotational speeds the valve train of an SI engine moves
from a quasi-static loading regime, in which the valve simply tracks
the cam profile, into an inertia-dominated regime in which the spring
force is the only mechanism that keeps the follower in contact with
the cam \citep{heywood1988,taylor1985}. Five interacting failure
modes emerge in this regime, and they are not independent: float can
trigger an impact that exceeds the residual coil-bind clearance,
and resonance can amplify both.

\begin{itemize}\setlength{\itemsep}{0.2em}
\item \emph{Valve float}, the spring force can no longer overcome
the inertial force on the cam deceleration ramp, the follower
separates from the cam, and re-contact occurs in an uncontrolled
manner \citep{heywood1988,taylor1985}.
\item \emph{Valve bounce}, impact-driven oscillation at seating
that degrades sealing and accelerates seat wear when the closing
velocity grows with cam acceleration \citep{heywood1988}.
\item \emph{Coil bind}, adjacent coils make metal-to-metal contact,
the spring stiffness rises abruptly, and the dynamic load is
transmitted through the coil stack \citep{shigley2020}.
\item \emph{Spring fatigue}, of the order of $10^{8}$ cycles per
200~h of operation under fluctuating shear loading between
seat-preload and full-lift stress; any surface or shot-peening
defect can dominate the failure process \citep{shigley2020}.
\item \emph{Resonance with cam harmonics}, coincidence of a cam
harmonic with the spring natural frequency drives surge, accelerates
coil-bind approach, and amplifies float \citep{heywood1988}.
\end{itemize}

The BowTie methodology, introduced next, provides a unified language
to describe how each of these modes contributes to a single top event.

% =================================================================
\subsection{The BowTie methodology}\label{sec:bg-bowtie}

The BowTie methodology organises a risk analysis around a single top
event: threats are arranged on the left side connected to the top
event through preventive barriers, and consequences are arranged on
the right side connected through mitigative barriers
\citep{deruijter2016bowtiereview}. Compared with FMEA, three
advantages are relevant here: (i)~the causal chain is made explicit
rather than compressed into a severity--occurrence--detectability
triplet; (ii)~each preventive barrier can be associated with a
measurable indicator and a pass criterion, which is the basis of the
quantitative coupling proposed in this paper; and (iii)~the
diagrammatic representation is accessible to designers, operators,
and reliability engineers simultaneously. The classical BowTie is
qualitative; quantitative extensions in the reliability literature
include Bayesian-network mapping
\citep{khakzad2012ress,khakzad2013bowtiebayesian}, Petri-net
prognostics \citep{vileiniskis2017ressbowtie}, adaptive
bow-tie formulations \citep{talebberrouane2021ressabt}, and
cloud-model and Dempster--Shafer evidential extensions
\citep{xie2021ressbowtie,sezer2023ressbowtie}, and full barrier
quantification under data scarcity
\citep{das2021comprehensive,wu2023novel}. Within
\emph{Quality and Reliability Engineering International} specifically,
recent work has extended the bow tie with operational
risk-influence factors \citep{hellas2025bobta}, improved the
resolution of FMEA/RPN prioritisation under linguistic uncertainty
\citep{tang2025pulfmea}, and applied single-component Monte-Carlo
propagation to structural reliability \citep{qian2023montecarlo}.
The QC-BowTie sits at the intersection of these three lines, barrier-based structuring, FMEA/RPN improvement, and Monte-Carlo
propagation, but differs in kind: its barrier states are derived
from \emph{deterministic, first-principles indicators} rather than
from event data or expert linguistic scores. Each preventive barrier
is mapped to an indicator from the analytical sizing of the component
with a literature-anchored pass criterion, and the indicators are
subsequently propagated through a Monte Carlo simulation to convert
the nominal-point status into manufacturing-realistic probabilities
of barrier compliance (Section~\ref{sec:res-mc}).

% =================================================================
\subsection{Related work}\label{sec:bg-related}

Several authors have addressed valve-train and engine-component
reliability from complementary perspectives. The work of
\citet{velezgodino2019overheadvalvetrain} on overhead valve trains in
urban buses is a representative case study published in
\textit{Engineering Failure Analysis}: the authors use field data and
metallurgical inspection to attribute the recurrent failure of an
in-service valve train to a combination of operational misuse and
inadequate inspection intervals, but they stop short of organising the
findings around a barrier-based diagram. \citet{soffritti2018wornvalvetrain}
report a four-cylinder diesel engine in which excessive wear at the
rocker-arm/pushrod and rocker-arm/valve interfaces appeared after only
1\,000~hours of operation; their root-cause analysis points to valve
misalignment with respect to the seat insert as the dominant trigger.
Both studies confirm that valve-train failures, once initiated,
propagate quickly into other components and motivate the need for a
preventive analysis framework.

Two further studies in \textit{Engineering Failure Analysis} address
elements of the valve-train system that are central to the present
work. \citet{xing2021valvespringEAC} investigate environment-assisted
cracking in high-strength valve springs and show that the fatigue limit
adopted at the design stage must be discounted to account for surface
condition and lubricant chemistry; this finding directly supports the
fatigue factor of safety adopted in
Section~\ref{sec:method-bowtie-quant}.
\citet{zhao2023camshaftV6} report lubrication failure in the camshaft
bearings of a V6 diesel engine, identifying an asperity-contact regime
that violates the design assumption of full hydrodynamic lubrication
and reduces the effective service life of the bearing pair; their
result reinforces the importance of the lubrication-related preventive
barriers shown on the left side of the BowTie diagram of
Section~\ref{sec:method-bowtie-qual}.
\citet{ko2022springfatigueflaws} extend the spring-fatigue discussion
beyond \textit{Engineering Failure Analysis}, showing in
\textit{Scientific Reports} that surface flaws as small as a few
micrometres on the oil-tempered wire of an automotive valve spring can
reduce the fatigue life by an order of magnitude; this result motivates
the conservative factor of safety adopted in
Table~\ref{tbl:quant-criteria}.

Table~\ref{tbl:related-cases} consolidates the four
\textit{Engineering Failure Analysis} case studies above and maps
their reported root causes to the preventive (P) and mitigative (M)
barriers and the deterministic indicators of the BowTie procedure
constructed in Section~\ref{sec:methodology}. The mapping is direct
for two of the four cases (Xing on environment-assisted cracking,
caught by the Goodman fatigue indicator; and the operational/
maintenance signature of V\'elez~Godi\~no, caught by the mitigative
side of the BowTie); is partial for the assembly-driven misalignment
of Soffritti, which the qualitative bow flags but which lies outside
the deterministic-indicator set adopted here; and is explicitly
\emph{outside} the present scope for the camshaft-bearing
lubrication failure of Zhao, which would require an additional
preventive barrier on the lubrication regime. The mixed mapping is
itself an outcome of the comparison: it identifies which classes of
failure the procedure already covers and which classes call for an
extension of the barrier set, and it is examined in detail in
Section~\ref{sec:res-litvalidation}.

\begin{table*}[!htbp]
\caption{Mapping of four documented valve-train and engine-component failure case studies from \textit{Engineering Failure Analysis} to the preventive (P) and mitigative (M) barriers and the deterministic indicators of the BowTie procedure proposed in this work. The notation P1--P5, M1--M5 refers to the barriers introduced in Section~\ref{sec:method-bowtie-qual}, and the indicators are those of Table~\ref{tbl:quant-criteria}. Reported root causes are paraphrased from the cited papers; no values are reproduced.}\label{tbl:related-cases}
\renewcommand{\arraystretch}{1.15}
\begin{tabular}{@{}p{0.19\textwidth}p{0.26\textwidth}p{0.26\textwidth}p{0.20\textwidth}@{}}
\toprule
Case study (component / engine) & Reported root cause & BowTie barrier triggered & Indicator that would flag it \\
\midrule
\citet{velezgodino2019overheadvalvetrain}; overhead valve train, urban bus & Operational misuse and inadequate inspection intervals & M3 (operating envelope at design) and M5 (predictive maintenance via vibration) & Mitigative side; outside the deterministic indicator set \\
\addlinespace
\citet{soffritti2018wornvalvetrain}; rocker-arm / pushrod interfaces, four-cylinder diesel ($\sim$1\,000~h) & Valve misalignment relative to the seat insert & Threat: out-of-tolerance assembly $\rightarrow$ candidate geometric/quality barrier P6 (extension) & Outside the present indicator set; flagged qualitatively by the bow chain \\
\addlinespace
\citet{xing2021valvespringEAC}; high-strength valve springs (premature fatigue) & Environment-assisted cracking due to surface condition and lubricant chemistry & P4 (high-fatigue-resistance spring material) & Goodman factor of safety $n_{f}$ with discounted shear endurance limit $S_{se}$ \\
\addlinespace
\citet{zhao2023camshaftV6}; camshaft bearings, V6 diesel & Asperity-contact regime violating the full-hydrodynamic-lubrication assumption & Threat: lubrication degradation; \emph{outside the present preventive set} & Outside scope; motivates a lubrication-regime barrier in procedure extensions \\
\bottomrule
\end{tabular}
\end{table*}

The methodological literature on the BowTie diagram itself has
matured in parallel. The review by \citet{deruijter2016bowtiereview}
remains the canonical reference for the qualitative formulation and
catalogues the variations adopted in the chemical, aerospace, and
nuclear industries. \citet{khakzad2013bowtiebayesian} propose a mapping
of the BowTie into a Bayesian network so that the static qualitative
diagram becomes dynamic and probabilistic; their construction is
adopted in part by the BowTie tool used in the offshore industry.
\citet{aust2019bowtieaircraft} apply the methodology to aircraft engine
borescope inspection and show that, even in the absence of
probabilistic data, the BowTie diagram is an effective communication
artefact between maintenance planners, inspectors, and reliability
engineers. The methodology continues to be extended in
\textit{Reliability Engineering \& System Safety} along two
complementary fronts:
\citet{xie2021ressbowtie} combine BowTie analysis with cloud-model
risk-matrix methods to evaluate fire and explosion risks in oil depots
where severity rubrics are inherently linguistic, and
\citet{sezer2023ressbowtie} couple the BowTie diagram with
Dempster--Shafer evidence theory and the HEART human-reliability
framework for cargo-tank crack risk analysis. These extensions
demonstrate the versatility of the underlying diagram when calibrated
with domain-specific quantitative or evidential layers.

It is also relevant to acknowledge a parallel branch of the literature
that extends the conventional FMEA in different directions: fuzzy-FMEA
formulations and analytic-hierarchy-process-based weighting (AHP-FMEA)
have been used in mechanical reliability to mitigate the loss of
resolution of the multiplicative Risk Priority Number (RPN). The
present work is complementary to those extensions: it does not aim
to refine the FMEA score but to replace the single-score logic with
a barrier-by-barrier deterministic check, and is therefore
positioned closer to the BowTie tradition than to fuzzy/AHP-FMEA.

To the best of the authors' knowledge, no published study applies the
BowTie methodology to the valve train of a high-rpm spark-ignition
cylinder head, and no study couples the BowTie diagram to the
analytical sizing of the intake valve and the valve spring as a
reproducible barrier-evaluation procedure. The present work fills both
gaps.
